% !PRINTED 391
\weberpage{391}
If the period of the continued fraction development of $\omega$ is
\weq{17}{[\alpha_0, \alpha_1 \ldots \alpha_{\nu-1}],}
then for $-1 : \omega'$ it is
\weq{18}{[\alpha_{\nu-1}, \alpha_{\nu-2} \ldots \alpha_0],}
that is, the reverse.

Now if $\omega$ is equivalent to $\omega'$, and hence also to $-1 : \omega'$, then
by 6. these two periods must agree with one another when one
begins at a suitable term, or the period of
$\omega$ must be \gesperrt{reversible}.

This means: in the continued fraction development
for $\omega$ two elements must admit of being so chosen that the development
reads the same whether one reads it proceeding from the one of these numbers to
the left or from the other to the right.
In symbols: for some suitably determined $k$ and
for $i = 0, 1, 2 \ldots \nu - 1$ we must have
\[
\alpha_i = \alpha_{\nu+k-i}
\]
If the period is reversible, then conversely
(by 6.) $\omega$ is equivalent to $-1 : \omega'$, and hence also to $\omega'$. From this,
then, the result:

8. \gesperrt{Two-sided numbers have in their continued
 fraction development a reversible period, and conversely
 numbers which have a reversible period
 in their continued fraction development
are two-sided.}

Finally we wish to remark further that by the property
of admitting development into a periodic continued fraction,
the real quadratic irrational numbers are distinguished from all other
numbers. For if $\omega$ is a number developed into a continued
fraction, and $\omega_n$, $\omega_m$ are two closing numbers
of this continued fraction, then
\[
\omega = \frac{P_n\omega_n + P_{n-1}}{Q_n\omega_n + Q_{n-1}} = \frac{P_m\omega_m + P_{m-1}}{Q_m\omega_m + Q_{m-1}}.
\]

Now if the continued fraction is periodic, no matter whether
the periodicity begins right from the start or only in the course
of the development, then for two different values of $n$
and $m$ we have $\omega_n = \omega_m$, and one thus obtains, by elimination of
%CONT>>
